% content/c00-theorie/en.tex — Chapter 0: What a physical theory supplies (English edition)
%
% Same structure, IDs and display equations as de.tex, which leads.  Migration of the first
% lecture (latex/lectures/en/lec01.tex is the English v0.2 text); see
% evidence/migrations/maps/vl01-map.json.  Voice: TJO's own sentences from the script and
% from his recorded lectures (transcripts/clean/v01.md, v02.md); timestamps are in de.tex.
% The transcripts are not a registered source and supply wording, examples and didactic
% remarks only, no physics.  Framing: docs/plan/BRIEF.md (author-stipulated, quoted
% verbatim here).  Physics of the prehistory: the author's lecture notes only
% (tpc-script-v0.2), scope "as stated in the author's lecture notes; no external source
% checked", or derived (evidence/derivations/c01/DERIVATION.md, part A).  Constants:
% CODATA 2022.  Order as in the lecture: prehistory and the two maxims first, then the
% framing of the course; section IDs are names, so c00.s4 precedes c00.s1.

\chapter{What a physical theory supplies}\label{c00}

We find ourselves in a quantum universe. It has always been so and always will be. Until
1900, however, we did not notice the quantum aspects of the universe, because classical
physics supplied us with a very good approximation to the physics of everyday life. That is
a great success and at the same time a problem: how is one to discover quantum aspects when
the classical approximation works so well? Quantum mechanics would be far easier to
understand if we could observe quantum aspects every day.

What stands in need of explanation is therefore not quantum mechanics, but the quality of
the classical approximation. One of the greatest challenges in learning quantum mechanics is
to understand why classical physics works so well. The keyword for that is
\emph{decoherence}; it belongs to advanced quantum mechanics, and we shall not manage it in
this course \srcref{tpc-script-v0.2}{lec01, section “The Birth of Quantum Mechanics”}.

% GAP-TJO: whether exercises on real quantum computers belong to the new course is TJO's
% decision (G-0070).
Today things are different from 1900. The first quantum computers have been built, and
through cloud services you can experiment yourself with real quantum computers. Quantum aspects are slowly becoming visible. Before we can play with them we need
the prehistory and the formal structure of the theory: on which foundations does it rest,
and which equations does one need?

First a remark on method, one that holds for the whole course. One can teach quantum
mechanics by working through the wrong, classical or semiclassical, models first, in order to
show afterwards that they do not work. That is not my style. From the very beginning we look
at everything correctly; nowhere do we work through the Bohr model of the atom. The
historical experiments I shall name nonetheless, but only as a list.

I am not recommending a textbook to you: in my opinion the best book on quantum mechanics
does not exist yet. There is no bad one, but equally no perfect one. Do read other books
alongside; the authoritative text is this one.

This chapter has two parts. First the findings on which classical physics failed, and the two
sentences in which I sum them up. Then the framing of the course: what a physical theory
supplies, and in which order we treat the theories. There is little calculation here; the
mathematics begins in Chapter~1 with probability.

\section{Why a new theory was needed}\label{c00.s4}

The prehistory of quantum mechanics is a series of observations which gradually called into
question the thesis that our universe is described by classical physics. As a line of proof
these old experiments hardly make sense today: we have long been convinced that quantum
mechanics is the right theory. They are good to know all the same. Whoever teaches in the
upper forms of a secondary school can tell an exciting story with them. Two of the findings
I give with their formula, the rest as a list.

\subsection{Black-body radiation}\label{c00.s4.1}

A black body reflects no incident radiation, yet at a temperature $\Temp$ it radiates all the
same, on account of its thermally excited atoms. When it is heated, its radiation has a
spectrum that depends only on the temperature and not on the material
\srcref{nobelprize-history-2026}{Max Planck, Facts}. A true black body is hard to make. The
usual approximation is a box with a small hole (Figure~\ref{c00.fig.hohlraum}) whose walls
all have the same temperature $\Temp$ and are in thermal equilibrium with the radiation
inside. We take this as a premise. Incident radiation vanishes inside through multiple
reflection, and approximately black-body radiation of temperature $\Temp$ emerges from the
hole.

Why does the incident radiation vanish? Suppose the wall sends back a fixed fraction $r$ of the
incident energy at every reflection, with $0 \leq r < 1$, and swallows the rest. After $n$
reflections the fraction $r^n$ of an entering ray is left, and
\begin{equation}\label{c00.eq.reflexion}\claim{c00.cl.reflexion}{derived}
  r^n \;\longrightarrow\; 0
  \qquad (n \to \infty) .
\end{equation}
A ray that finds the small hole again only after many reflections carries hardly any energy
out. The hole therefore acts approximately like a body that swallows everything. About the
radiation that comes out of the hole this calculation says nothing; for that one needs the
premise of thermal equilibrium.

\begin{tpcfigure}{c00.fig.hohlraum}
  \demo{c00.demo.hohlraum}
  \caption{A black body as a box with a small hole. Drawn are the first reflections of an
    entering ray according to the law of reflection; none of them hits the hole. If the
    wall keeps the fraction $1 - r$ of the energy at every reflection, only the fraction
    $r^n$ is left after $n$ reflections, equation~\eqref{c00.eq.reflexion}.}
\end{tpcfigure}

What was shocking around 1900 was that the spectrum of this radiation cannot be explained
with classical physics. The right spectrum follows from the assumption that light of
frequency $\nu$ consists of light quanta of energy
\begin{equation}\label{c00.eq.lichtquant}\claim{c00.cl.lichtquant}{source-backed}
  E = \hPlanck \nu
\end{equation}
The constant $\hPlanck$ is called Planck's constant. Planck originally postulated it as an
“auxiliary quantity” in deriving the radiation law. The law gives how much energy per volume
falls on the frequencies between $\nu$ and $\nu + \dd\nu$:
\begin{equation}\label{c00.eq.planck}\claim{c00.cl.planck}{source-backed}
  u(\nu, \Temp)\,\dd\nu
  = \frac{8\pi \hPlanck \nu^3}{\clight^3}\,\frac{\dd\nu}{e^{\hPlanck\nu/(\kB \Temp)} - 1} .
\end{equation}
Here $\clight$ is the speed of light and $\kB$ Boltzmann's constant. We shall derive this law
ourselves in the statistical part of the course, from the assumption that light is
quantised. Here it is only used \srcref{tpc-script-v0.2}{lec01, prehistory, item 1}.

Theoretical physics rarely uses $\hPlanck$ today; predominantly one uses the \emph{reduced
Planck constant}
\begin{equation}\label{c00.eq.hbar}\claim{c00.cl.hbar}{source-backed}
  \hbar = \frac{\hPlanck}{2\pi} .
\end{equation}
The numerical values are
$\hPlanck = \num[c00.cl.h-wert]{6.62607015 \cdot 10^{-34}}{J\,s}$ and
$\hbar = \num[c00.cl.hbar-wert]{1.054571817 \cdot 10^{-34}}{J\,s}$; the second is
truncated to ten digits.

What does the constant $\hPlanck$ do in the radiation law? For $\nu > 0$ and $\Temp > 0$ we
look at the exponent $x = \hPlanck\nu/(\kB \Temp)$. For small $x$, $e^x - 1$ is nearly equal
to $x$. Inserting this into \eqref{c00.eq.planck}, $\hPlanck$ cancels and the function
\begin{equation}\label{c00.eq.grenzform}\claim{c00.cl.grenzform}{derived}
  u_{\mathrm{kl}}(\nu, \Temp) = \frac{8\pi \nu^2 \kB \Temp}{\clight^3},
  \qquad
  \frac{u(\nu, \Temp)}{u_{\mathrm{kl}}(\nu, \Temp)} = \frac{x}{e^{x} - 1}
  \quad \text{\tpcifenglish{with}{mit}}\ x = \frac{\hPlanck\nu}{\kB \Temp} .
\end{equation}
remains. The second equation holds exactly for $\nu > 0$ and $\Temp > 0$; one obtains it by
dividing \eqref{c00.eq.planck} by $u_{\mathrm{kl}}$. The ratio $x/(e^x - 1)$ is close to $1$
for small $x$ and falls to $0$ for large $x$\zoomref{c00.v1}. The limiting form
$u_{\mathrm{kl}}$ no longer contains $\hPlanck$.
% GAP-TJO: that u_kl is the prediction of classical physics (Rayleigh-Jeans, "ultraviolet
% catastrophe") rests only on the old figure file, not on a registered source (G-0024).

\begin{tpcfigure}{c00.fig.planck}
  \demo{c00.demo.planck}
  \caption{The radiation law~\eqref{c00.eq.planck} and its limiting form
    \eqref{c00.eq.grenzform}. (a) The energy density $u$ for three temperatures $T_0$,
    $\tfrac56 T_0$ and $\tfrac23 T_0$ (solid, dashed, dash-dotted), plotted against
    $\hPlanck\nu/(\kB T_0)$ in units of $u_0 = 8\pi(\kB T_0)^3/(\hPlanck^2\clight^3)$. The
    thin dashed parabola is the limiting form $u_{\mathrm{kl}}$ for the highest temperature
    $T_0$. Notice how at small frequencies it lies approximately on the solid curve to which
    it belongs, and only then parts company with it; after that it runs off the top of the
    picture. (b) The ratio $u/u_{\mathrm{kl}} = x/(e^x-1)$ for $\nu > 0$ and $\Temp > 0$. It
    depends only on $x = \hPlanck\nu/(\kB\Temp)$. The hatched area is the distance from the
    limiting form. It becomes noticeable as soon as $\hPlanck\nu$ reaches the size of
    $\kB\Temp$.}
\end{tpcfigure}

Figure~\ref{c00.fig.planck} shows both. The limiting form grows like $\nu^2$ without end,
while the curves of the radiation law have a maximum and fall off afterwards. The difference
becomes visible as soon as the energy $\hPlanck\nu$ of a light quantum is no longer small
compared with $\kB\Temp$.

\begin{vertiefung}[id=c00.v1,
    titel={The limiting form, precisely},
    wozu={The Core says that the ratio $x/(e^x-1)$ is close to $1$ for small $x$ and falls
      to $0$ for large $x$. Here are two bounds that prove it, and a consequence: the
      limiting form predicts an infinite energy density, the radiation law a finite one.},
    voraussetzung={Core up to equation~\eqref{c00.eq.grenzform}; the exponential series and
      integrals of powers. The rest of the chapter does not need this block.},
    zurueck={For $\hPlanck\nu$ small compared with $\kB\Temp$ the radiation law and the
      limiting form agree approximately, their ratio tends to $1$; for large frequencies the
      limiting form is wrong, and arbitrarily wrong.}]
  For $x > 0$, $e^x = 1 + x + \tfrac12 x^2 + \tfrac16 x^3 + \dots$ is a sum of positive
  terms. Two estimates for $e^x - 1$ follow. Dropping all terms from the third on makes the
  sum smaller: $e^x - 1 \geq x + \tfrac12 x^2$. Factoring out $x$ instead leaves
  $1 + \tfrac12 x + \tfrac16 x^2 + \dots$ in the bracket, and each of its terms is at most as
  large as the corresponding term $1 + x + \tfrac12 x^2 + \dots$ of the series for $e^x$:
  so $e^x - 1 \leq x\,e^x$. Together, for all $x > 0$,
  \begin{equation}\label{c00.eq.schranken}\claim{c00.cl.schranken}{derived}
    e^{-x} \;\leq\; \frac{x}{e^{x} - 1} \;\leq\; \frac{1}{1 + \tfrac12 x} .
  \end{equation}
  As $x \to 0$ both bounds tend to $1$, as $x \to \infty$ both tend to $0$. At $x = 1$, that
  is $\hPlanck\nu = \kB\Temp$, the ratio lies between $1/e$ and $\tfrac23$.

  Now to the total energy density, the integral over all frequencies. For the limiting form,
  \begin{equation}\label{c00.eq.divergenz}\claim{c00.cl.divergenz}{derived}
    \int_0^{\nu_1} u_{\mathrm{kl}}(\nu, \Temp)\,\dd\nu
    = \frac{8\pi \kB \Temp}{3\,\clight^3}\,\nu_1^3 ,
  \end{equation}
  and this grows beyond every bound with the upper limit $\nu_1$. The limiting form predicts
  an infinite energy density for every temperature $\Temp > 0$.

  For the radiation law the substitution $\nu = x\,\kB\Temp/\hPlanck$ gives
  \begin{equation}\label{c00.eq.gesamt}\claim{c00.cl.gesamt}{derived}
    \begin{aligned}
      \int_0^{\infty} u(\nu, \Temp)\,\dd\nu
      &= \frac{8\pi (\kB\Temp)^4}{\hPlanck^3 \clight^3} \int_0^\infty \frac{x^3}{e^x - 1}\,\dd x , \\
      \int_0^\infty \frac{x^3}{e^x - 1}\,\dd x &\leq \frac13 + 120 .
    \end{aligned}
  \end{equation}
  The bound on the integral comes in two pieces. For $0 < x \leq 1$,
  \eqref{c00.eq.schranken} gives $x^3/(e^x-1) \leq x^2$, and
  $\int_0^1 x^2\,\dd x = \tfrac13$. For $x \geq 1$ one uses the term $x^5/120$ of the
  exponential series: $e^x - 1 \geq x^5/120$, so $x^3/(e^x-1) \leq 120/x^2$, and
  $\int_1^\infty 120\,x^{-2}\,\dd x = 120$. The integral is a pure number. The total energy
  density of the radiation law is therefore finite and proportional to $\Temp^4$. We need the
  value of the number only in the statistical part.
\end{vertiefung}

\subsection{Photoelectric effect}\label{c00.s4.2}

The second experiment concerns light and matter. Light can release electrons, the
photoelectrons. The finding: for a fixed material the \emph{maximum} kinetic energy of the
photoelectrons depends on the frequency of the light, not on its intensity. Classically that
is very strange. Everywhere else more intensity means more energy; here the energy of the
fastest electrons depends on the frequency alone. In the one-photon model the balance reads
\begin{equation}\label{c00.eq.photoeffekt}\claim{c00.cl.photoeffekt}{source-backed}
  E_{\mathrm{kin,max}} = \tfrac12 \massE v_{\max}^2 = \hPlanck\nu - W
\end{equation}
with the mass $\massE$ of the electron, its largest occurring speed $v_{\max}$ and the
binding energy $W$ of the material \srcref{tpc-script-v0.2}{lec01, prehistory, item 2}. The
kinetic energy is written here non-relativistically as $\tfrac12 \massE v^2$; this assumes
that the electrons are slow compared with light. This observation pointed to light quanta
with the energy~\eqref{c00.eq.lichtquant}.

Let us check the balance in the one-photon model. A light quantum gives its energy
$\hPlanck\nu$ to a single electron. To get out, this electron has to spend an energy $W'$;
$W$ is the smallest such energy, so $W' \geq W$. For an electron that actually gets out,
a kinetic energy is left, and that is never negative:
\begin{equation}\label{c00.eq.schwelle}\claim{c00.cl.schwelle}{derived}
  0 \;\leq\; \tfrac12 \massE v^2 = \hPlanck\nu - W' \;\leq\; \hPlanck\nu - W,
  \qquad \text{\tpcifenglish{hence}{also}} \qquad
  \nu \geq \frac{W}{\hPlanck} .
\end{equation}
Equation~\eqref{c00.eq.photoeffekt} thus gives the largest kinetic energy, the one with
$W' = W$. Below the threshold frequency $W/\hPlanck$ a single light quantum releases no
electron in this model. For fixed $W$, $E_{\mathrm{kin,max}}$ is a straight line in $\nu$ with slope $\hPlanck$, starting at $\nu = W/\hPlanck$
(Figure~\ref{c00.fig.photoeffekt}). The slope is the same for every value of $W$.

\begin{tpcfigure}{c00.fig.photoeffekt}
  \demo{c00.demo.photoeffekt}
  \caption{The maximum kinetic energy~\eqref{c00.eq.photoeffekt} of the photoelectrons,
    plotted for two fixed values of the binding energy, $W_1$ (solid) and
    $W_2 = \tfrac85 W_1$ (dashed; the ratio is chosen for the picture). Both lines have the
    same slope $\hPlanck$: a frequency difference $\nu' - \nu$ goes with the energy
    difference $\hPlanck(\nu' - \nu)$. Each line starts at its threshold
    $\nu = W/\hPlanck$ from equation~\eqref{c00.eq.schwelle}. The dotted continuations meet
    the axis at $-W$; there the kinetic energy would be negative.}
\end{tpcfigure}

\subsection{Further findings}\label{c00.s4.3}

\begin{enumerate}
  \item \emph{Stability of atoms, line spectra, Franck-Hertz experiment.} This is really an
    observation: atoms are stable, and classically that cannot be explained. By the
    Liénard-Wiechert potentials, which you know from electrodynamics, every accelerated point
    charge radiates. An electron on an orbit around the nucleus would therefore lose energy
    all the time and ought to crash into the nucleus, which it does not do. Line spectra and
    the experiments of Franck and Hertz indicate that only certain discrete energy levels
    are allowed in atoms. Only quantum mechanics can explain this; we shall discuss the line
    spectrum of hydrogen in detail in Chapter~11.
  \item \emph{Compton effect.} Light is scattered off electrons. How energy and momentum are
    shared between light quantum and electron is described correctly with light quanta of
    energy~\eqref{c00.eq.lichtquant} and momentum $p = \hPlanck/\lambda$ from
    equation~\eqref{c00.eq.impuls}, together with conservation of energy and momentum.
  \item \emph{Stern-Gerlach experiment.} This experiment matters particularly to me. A beam
    of silver atoms passes through an inhomogeneous magnetic field. Classically it should
    spread out into a continuous distribution; instead two separate beams are observed.
    Angular momentum is quantised: at fixed angular-momentum quantum number $\lqn$ its
    component in the direction of the field takes only finitely many values (space
    quantisation). From such quantised angular momenta one can build qubits today,
    the building blocks of quantum computers.
  \item \emph{Interference of electrons.} The diffraction of electrons from crystals
    (Davisson-Germer experiment) gave an indication of wave properties of the electron.
    Light has wave but also particle properties; the photoelectric and Compton effects were
    about that. Electrons have particle but also wave properties.
\end{enumerate}
The momentum of a light quantum of wavelength $\lambda$ is
\begin{equation}\label{c00.eq.impuls}\claim{c00.cl.impuls}{source-backed}
  p = \frac{\hPlanck}{\lambda}
\end{equation}
\srcref{tpc-script-v0.2}{lec01, prehistory, items 3 to 6}. The history with names and dates is
good to know; you will not need it for the rest of the course\zoomref{c00.v2}.

\begin{vertiefung}[id=c00.v2,
    titel={Names and dates},
    wozu={The Core names the findings without their history. Here are the dates and names,
      for those who want to tell the story.},
    voraussetzung={Section~\ref{c00.s4}. The rest of the course does not need this block.},
    zurueck={The mathematical form of quantum mechanics and its statistical interpretation
      arose in the years 1925 to 1927.}]
  % Dates corrected after REVIEW-lec1 (M-0050); Lenard's year is not established by the
  % registered pages and is left out (G-0025).
  The text of this block follows my lecture notes \srcref{tpc-script-v0.2}{lec01,
  prehistory}; the years of the experiments have been checked against the laureate pages of
  the Nobel Foundation \srcref{nobelprize-history-2026}{Facts pages on Planck, Franck,
  Stern, de Broglie, Einstein, Davisson; Crossref record of Phys. Rev. 30, 705}. One can argue
  well:
  \begin{equation}\label{c00.eq.geburtsjahr}\claim{c00.cl.geschichte}{source-backed}
    \text{\tpcifenglish{The year of birth of quantum mechanics is 1926.}{Das Geburtsjahr der Quantenmechanik ist 1926.}}
  \end{equation}

  \begin{tpctable}{c00.tab.jahre}
    \caption{Dates and names for the findings of Section~\ref{c00.s4}. The years follow the
      laureate pages of the Nobel Foundation and, for Davisson and Germer, the year their
      paper appeared in the Physical Review.}
    \begin{tabular}{ll}
      \toprule
      Finding & Names, year \\
      \midrule
      Radiation law, light quanta & Planck, 1900 \\
      Photoelectric effect & Lenard; Einstein, from 1905 \\
      Franck-Hertz experiment & Franck and Hertz, 1914 \\
      Stern-Gerlach experiment & Stern and Gerlach, 1922 \\
      Wave properties of the electron & de Broglie, 1924 (postulate) \\
      Davisson-Germer experiment & Davisson and Germer, 1927 \\
      \bottomrule
    \end{tabular}
  \end{tpctable}

  In 1926 Schrödinger published an exciting series of papers on wave mechanics. A parallel
  approach began in 1925 with reflections by Heisenberg; from it Born and Jordan with
  Heisenberg developed matrix mechanics. So one had, as it were, two quantum theories, wave
  mechanics and matrix mechanics. Still in 1926 Schrödinger showed that wave mechanics and
  matrix mechanics are equivalent. Heisenberg's work also led Dirac to a version of quantum
  mechanics of his own. Born supplied, likewise in 1926, the statistical interpretation of
  the theory. In my view Born's contribution is the
  most important aspect of the whole theory.

  Almost at the same time the theory took on the mathematical form it has today. Born took
  advice from Norbert Wiener on turning matrices into operators. Hilbert gave a lecture
  course on quantum mechanics in Göttingen in the winter semester 1926/27, and alongside it
  his young assistant John von Neumann introduced the Hilbert space, density matrices and the
  entropy of quantum states. Physicists did not read it at first. Wigner and Weyl clarified
  the role of group theory. It is remarkable that, in spite of the theory's fantastic
  successes, its interpretation remained disputed; to this day there is no contradiction
  with quantum mechanics.

  \begin{historical}
    You are, by the way, in luck if you read German: the original papers of Planck,
    Schrödinger, Heisenberg and Born are written in German. I did not have that opportunity
    when I learnt quantum mechanics.
  \end{historical}
\end{vertiefung}

\subsection{Two maxims}\label{c00.s4.4}

All these observations together are indications that quantum mechanics is the right theory
of the universe. I sum them up in two sentences
\srcref{tpc-script-v0.2}{lec01, the two yellow boxes}. The first:
\begin{equation}\label{c00.eq.welle-teilchen}\claim{c00.cl.welle-teilchen}{source-quoted}
  \begin{aligned}
    &\text{\tpcifenglish{Matter and radiation possess}{Materie und Strahlung besitzen}} \\
    &\text{\tpcifenglish{both wave and particle properties.}{sowohl Wellen- als auch Teilcheneigenschaften.}}
  \end{aligned}
\end{equation}
The second is the main message of this course:
\begin{equation}\label{c00.eq.kernsatz}\claim{c00.cl.kernsatz}{source-quoted}
  \begin{aligned}
    &\text{\tpcifenglish{Quantum mechanics is a probabilistic theory.}{Die Quantenmechanik ist eine probabilistische Theorie.}} \\
    &\text{\tpcifenglish{It therefore gives us predictions}{Sie liefert uns somit nur Vorhersagen}} \\
    &\text{\tpcifenglish{only in the form of probabilities.}{in Form von Wahrscheinlichkeiten.}}
  \end{aligned}
\end{equation}
This second statement stands above the first. “Wave and particle properties” remains a
little vague; the probabilistic statement is precise and underlies every quantum phenomenon.
Classically the predictions are definite: the particle is at this place at this time.
Probabilities enter classical physics purely through missing information, for instance in
statistical physics; there are no fundamental probabilities there. In quantum mechanics it is
the other way round: it only says with which probability an event takes place. This
statistical interpretation was supplied by Born, and for me it is the most important aspect
of the theory.

\begin{merksatz}[The most important statement about quantum mechanics]
  Quantum mechanics is a probabilistic theory: it gives predictions only in the form of
  probabilities. Please take this sentence away with you. In every oral examination I ask
  what the most important statement about quantum mechanics is.
\end{merksatz}

I want to be honest: the interpretation of quantum mechanics is still disputed today, despite
many ideas. To my taste the probabilistic statement is the right one. Quantum mechanics is the
best model we have in physics, and to this day there is no contradiction with it.

We now leave the exciting history of quantum mechanics behind. From now on:
\begin{equation}\label{c00.eq.annahme}\claim{c00.cl.annahme}{source-quoted}
  \begin{aligned}
    &\text{\tpcifenglish{Our assumption is that quantum mechanics}{Unsere Annahme ist, dass die Quantenmechanik}} \\
    &\text{\tpcifenglish{is the correct description of the universe.}{die korrekte Beschreibung des Universums ist.}}
  \end{aligned}
\end{equation}

\section{The framing of the course}\label{c00.s1}

% The three sentences of this section are docs/plan/BRIEF.md:50, 44 and 45 verbatim
% (without the labels).  GAP-TJO: the German wording in de.tex is a translation (G-0022);
% reading P0 as the second maxim said for every theory compares the two wordings (G-0023 c).
From the second maxim I make a principle for the whole course. It holds not only for quantum
mechanics but for every physical theory:
\begin{equation}\label{c00.eq.p0}\claim{c00.cl.p0}{author-stipulated}
  \begin{aligned}
    &\text{\tpcifenglish{Physical theories are probabilistic theories:}{Physikalische Theorien sind probabilistische Theorien:}} \\
    &\text{\tpcifenglish{predictions in the form of probabilities.}{Vorhersagen in Form von Wahrscheinlichkeiten.}}
  \end{aligned}
\end{equation}
I call this sentence P0. It is the first of the postulates that we set up step by step; with
each further postulate the theory becomes more special, until it is quantum mechanics. Two
further stipulations belong with it:
\begin{equation}\label{c00.eq.liefert}\claim{c00.cl.liefert}{author-stipulated}
  \begin{aligned}
    &\text{\tpcifenglish{A physical theory supplies data for}{Eine physikalische Theorie liefert Daten für}} \\
    &\text{\tpcifenglish{preparations + dynamics + measurement.}{Präparationen + Dynamik + Messung.}}
  \end{aligned}
\end{equation}
\begin{equation}\label{c00.eq.probabilistisch}\claim{c00.cl.probabilistisch}{author-stipulated}
  \begin{aligned}
    &\text{\tpcifenglish{All physical theories are}{Alle physikalischen Theorien sind}} \\
    &\text{\tpcifenglish{probabilistic theories (I think of deterministic}{probabilistische Theorien (deterministische fasse ich}} \\
    &\text{\tpcifenglish{as a subset of probabilistic).}{als Teilmenge der probabilistischen auf).}}
  \end{aligned}
\end{equation}
Sentence~\eqref{c00.eq.liefert} names three ingredients; the next section shows where
preparation and measurement occur in an experiment. Sentence~\eqref{c00.eq.probabilistisch} says that the definite
predictions of classical physics fit into this frame too. What that looks like in the
simplest model we see in Chapter~1: there are preparations for which every measurement gives
a fixed value with probability $1$.
% GAP-TJO: a definition of deterministic theories including dynamics is outstanding
% (G-0023 d).

\section{Every experiment in three steps}\label{c00.s2}

To the theorist all experiments in physics look the same
\srcref{tpc-script-v0.2}{lec01, box “Every experiment in three steps”}:
\begin{equation}\label{c00.eq.drei-stadien}\claim{c00.cl.drei-stadien}{source-quoted}
  \begin{aligned}
    &\text{\tpcifenglish{To the theorist all experiments in physics}{Für den Theoretiker sehen alle Experimente}} \\
    &\text{\tpcifenglish{look the same; they consist of three stages.}{der Physik gleich aus; sie bestehen aus drei Stadien.}}
  \end{aligned}
\end{equation}
\begin{enumerate}
  \item \emph{Preparation.} Something is set up: the die, in the laboratory for instance a
    beam of light.
  \item \emph{Run.} The experiment runs; we do something with what was prepared.
  \item \emph{Measurement.} At the end we look and obtain one of the possible outcomes.
\end{enumerate}
Figure~\ref{c00.fig.experiment} shows the scheme. An experimental physicist would not accept
this simplification; abstractly, though, every experiment looks like this: prepare, run,
measure.

\begin{tpcfigure}{c00.fig.experiment}
  \demo{c00.demo.experiment}
  \caption{The three stages of an experiment. The brace combines the first two into a single
    preparation. The preparation is drawn solid, the run with a heavier stroke, the
    measurement dashed.}
\end{tpcfigure}

The first two steps can be combined into \emph{one} more complicated preparation. A loaded
die is nothing but an ordinary die to which something was done beforehand. Where one draws
the line is a matter of taste; that is why one can often make do with two notions,
preparation and measurement. These two also occur in sentence~\eqref{c00.eq.liefert}.
% GAP-TJO: that the script's "run" is the brief's "dynamics" is stated in neither source
% (G-0023 a).

What a preparation is, precisely, we fix provisionally with a definition adapted from
Plávala's presentation:

\tpcdef{D1}

The yes-no questions of the definition are the simplest measurements: those with only two
possible outcomes, such as a detector that says click or no click. In Chapter~1 they are the
indicator functions, and Chapter~2 defines them in general.
% GAP-TJO: D1 is adapted from Plávala; TJO's deviations from Ludwig-Werner are unwritten
% (G-0023 e).

\section{From small to large}\label{c00.s3}

The order of the material is fixed too:
\begin{equation}\label{c00.eq.leiter}\claim{c00.cl.leiter}{author-stipulated}
  \begin{aligned}
    &\text{\tpcifenglish{Our entire presentation works like this:}{Unsere ganze Darstellung geht so vor:}} \\
    &\text{\tpcifenglish{we go from microscopic to macroscopic.}{Wir gehen vom Mikroskopischen zum Makroskopischen.}}
  \end{aligned}
\end{equation}
Figure~\ref{c00.fig.leiter} shows the stages. From stage to stage one “zooms out”; at one
point information is lost explicitly.

\begin{tpcfigure}{c00.fig.leiter}
  \demo{c00.demo.leiter}
  \caption{The ladder of theories of this course, from the bottom (microscopic) to the top
    (macroscopic). The labels of the six stages and the five arrows are selected excerpts
    from the stipulation to which sentence~\eqref{c00.eq.leiter} refers. There the last stage
    reads in full “thermodynamics as the ultimate effective theory of all physics at the
    macroscale”. GPT stands for “general probabilistic theory”, QM for quantum mechanics.
    The course begins at the bottom: Chapters~1 and~2 treat stage~1.}
\end{tpcfigure}

Every stage comes with the same thought:
\begin{equation}\label{c00.eq.effektiv}\claim{c00.cl.effektiv}{author-stipulated}
  \begin{aligned}
    &\text{\tpcifenglish{At each stage the student learns that}{Auf jeder Stufe lernt man, dass}} \\
    &\text{\tpcifenglish{every theory is an effective theory of the one below it:}{jede Theorie eine effektive Theorie der darunterliegenden ist:}} \\
    &\text{\tpcifenglish{an approximation to the fundamental underlying physics.}{eine Näherung an die grundlegende Physik darunter.}}
  \end{aligned}
\end{equation}
What “zooming out” means at a particular step is the subject of the chapter that takes that
step. Here only the order is given.

Our next goal is to build the mathematical foundations of quantum mechanics. The numbers it
delivers as results are probabilities. To reach the goal we must therefore first treat
probability mathematically. That is what the next chapter does.
